\par\medskip
\SourceLeaf{072}\pdfbookmark[2]{54.}{article-054}
\begin{center}54.\end{center}
Relationes motus ad eclipticam in praecc. explicatae manifesto perinde valebunt, etiamsi pro ecliptica quodvis aliud planum substituatur, si modo situs plani orbitae ad hoc planum innotuerit; expressiones longitudo et latitudo autem tunc supprimendae erunt. Offert itaque se problema: \textit{e situ cognito plani orbitae aliusque plani novi ad eclipticam derivare situm plani orbitae ad hoc novum planum.} Sint $n\AscendingNode$, $\AscendingNode\AscendingNode'$, $n\AscendingNode'$ partes circulorum maximorum, quos planum eclipticae, planum orbitae planumque novum in sphaera coelesti proiiciunt (Fig. 2). Ut inclinatio circuli secundi ad tertium locusque nodi ascendentis absque ambiguitate assignari possit, in circulo tertio alterutra directio eligi debebit tamquam ei analoga, quae in ecliptica est secundum ordinem signorum; sit haec in fig. nostra directio ab $n$ versus $\AscendingNode'$. Praeterea duorum hemisphaeriorum, quae circulus $n\AscendingNode'$ separat, alterum censere oportebit analogum haemisphaerio boreali, alterum australi: haec vero hemisphaeria sponte iam sunt distincta, quatenus id semper quasi boreale spectatur, quod in circulo secundum ordinem signorum progredienti\footnote{Puta in \textit{interiori} sphaerae superficie, quam figura nostra repraesentat.} a dextra est. In figura igitur nostra sunt $\AscendingNode$, $n$, $\AscendingNode'$ nodi ascendentes circuli secundi in primo, tertii in primo, secundi in tertio; $180^0-n\AscendingNode\AscendingNode'$, $\AscendingNode n\AscendingNode'$, $n\AscendingNode'\AscendingNode$ inclinationes secundi ad primum, tertii ad primum, secundi ad tertium. Pendet itaque problema nostrum a solutione trianguli sphaerici, ubi e latere uno angulisque adiacentibus reliqua sunt deducenda. Praecepta vulgaria, quae in trigonometria sphaerica pro hoc casu traduntur, tamquam abunde nota supprimimus: commodius autem methodus alia in usum vocatur ex aequationibus quibusdam petita, quae in libris nostris trigonometricis frustra quaeruntur. Ecce has aequationes, quibus in sequentibus frequenter utemur: designant $a$, $b$, $c$ latera trianguli sphaerici atque $A$, $B$, $C$ angulos illis resp. oppositos:\HistoricalNoteDouble{Handschriftliche Bemerkung:] Si duo triangula communia habent duo latera $b$, $c$, erit
\[
\begin{aligned}
\mathrm{I.}\quad&\frac{\sin\frac{1}{2}(a'-a)}{\sin\frac{1}{2}(A'-A)}=\frac{\sin b\sin\frac{1}{2}(C'+C)}{\cos\frac{1}{2}(B'-B)}=\frac{\sin c\sin\frac{1}{2}(B'+B)}{\cos\frac{1}{2}(C'-C)}\\[.5em]
\mathrm{II.}\quad&-\frac{\sin\frac{1}{2}(a'+a)}{\sin\frac{1}{2}(A'-A)}=\frac{\sin b\cos\frac{1}{2}(C'+C)}{\sin\frac{1}{2}(B'-B)}=\frac{\sin c\cos\frac{1}{2}(B'+B)}{\sin\frac{1}{2}(C'-C)}\\[.5em]
\mathrm{III.}\quad&-\frac{\sin\frac{1}{2}(a'-a)}{\sin\frac{1}{2}(A'+A)}=\frac{\sin b\sin\frac{1}{2}(C'-C)}{\cos\frac{1}{2}(B'+B)}=\frac{\sin c\sin\frac{1}{2}(B'-B)}{\cos\frac{1}{2}(C'+C)}\\[.5em]
\mathrm{IV.}\quad&\frac{\sin\frac{1}{2}(a'+a)}{\sin\frac{1}{2}(A'+A)}=\frac{\sin b\cos\frac{1}{2}(C'-C)}{\sin\frac{1}{2}(B'+B)}=\frac{\sin c\cos\frac{1}{2}(B'-B)}{\sin\frac{1}{2}(C'+C)}\cdot
\end{aligned}
\]
}
\SourceLeaf{073}
\[
\mathrm{I.}\quad\frac{\sin\frac{1}{2}(b-c)}{\sin\frac{1}{2}a{,}}=\frac{\sin\frac{1}{2}(B-C)}{\cos\frac{1}{2}A}
\]
\[
\mathrm{II.}\quad\frac{\sin\frac{1}{2}(b+c)}{\sin\frac{1}{2}a}=\frac{\cos\frac{1}{2}(B-C)}{\sin\frac{1}{2}A}
\]
\[
\mathrm{III.}\quad\frac{\cos\frac{1}{2}(b-c)}{\cos\frac{1}{2}a}=\frac{\sin\frac{1}{2}(B+C)}{\cos\frac{1}{2}A}
\]
\[
\mathrm{IV.}\quad\frac{\cos\frac{1}{2}(b+c)}{\cos\frac{1}{2}a}=\frac{\cos\frac{1}{2}(B+C)}{\sin\frac{1}{2}A}\cdot
\]
Quamquam demonstrationem harum propositionum brevitatis caussa hic praeterire oporteat, quisque tamen earum veritatem in triangulis, quorum nec latera nec anguli $180^0$ excedunt, haud difficile confirmare poterit. Quodsi quidem idea trianguli sphaerici in maxima generalitate concipitur, ut nec latera nec anguli ullis limitibus restringantur (quod plurima commoda insignia praestat, attamen quibusdam dilucidationibus praeliminaribus indiget), casus existere possunt, ubi in cunctis aequationibus praecedentibus signum mutare oporteat; quoniam vero signa priora manifesto restituuntur, simulac unus angulorum vel unum laterum $360^0$ augetur vel diminuitur, signa, qualia tradidimus, semper tuto retinere licebit, sive e latere angulisque adiacentibus reliqua determinanda sint, sive ex angulo lateribusque adiacentibus; semper enim vel quaesitorum valores ipsi vel $360^0$ a veris diversi hisque adeo aequivalentes per formulas nostras elicientur. Dilucidationem copiosiorem huius argumenti ad aliam occasionem nobis reservamus: quod vero praecepta, quae tum pro solutione problematis nostri tum in aliis occasionibus formulis istis superstruemus, in omnibus casibus generaliter valent, tantisper adiumento inductionis rigorosae, i. e. completae omnium casuum enumerationis, haud difficile comprobari poterit.
